\section{Introduction}

Magnetic damping, also known as Gilbert damping ($\alpha$) in the context of the Landau-Lifshitz-Gilbert (LLG) equation \cite{eriksson2017atomistic}, plays a crucial role in determining the rate at which energy and angular momentum dissipate in a magnetic system. It thus generally dictates the timescale of processes involving a magnetic transition from an excited state to an equilibrium state, magnetic domain wall motion, and spin-wave propagation. These are the fundamental processes that govern the performance of spintronics \cite{chumak2015magnon} and magnonics \cite{barman20212021} devices. Reduced magnetic damping can lead to higher efficiency and faster operations in these devices by minimizing energy dissipation during their operation. Hence, a detailed investigation into the fundamental physics involved in the mechanism of damping is of great interest \cite{schoen2016ultra, jungwirth2016antiferromagnetic, chumak2015magnon, yang2022two}. 
Previous theoretical and experimental research has established that $\alpha$ is not a simple scalar quantity (as is usually assumed), but instead a tensorial quantity \cite{PhysRevB.72.064450, PhysRevB.78.020410, Mondal2018, PhysRevB.98.104406, PhysRevLett.108.057204, PhysRevLett.102.086601} that is temperature-dependent \cite{Zhao2016, PhysRevB.87.014430, Hiramatsu2021}, non-local in real-space \cite{thonig2014gilbert, PhysRevMaterials.2.013801, PhysRevB.72.064450, PhysRevB.78.020410, PhysRevB.105.155151, PhysRevB.108.014433, PhysRevLett.113.237204, Weissenhofer2024, Reyesosorio2023}, and presents an anisotropic behavior -- as outlined in various studies  \cite{PhysRevMaterials.2.013801, PhysRevB.103.L220405, PhysRevLett.122.117203, brinker2022generalization, PhysRevB.104.024404}.
Therefore, fundamental understanding and materials engineering at the local (atomistic) scale compose, in principle, a possible route to optimum properties with desired damping magnitude.

Apart from their exceptional properties, including high Curie temperature, high saturation magnetization and high permeability \cite{sundar2005soft}, Fe-Co alloys constitute also a suitable platform to the fundamental investigation of the magnetic damping. It was experimentally demonstrated that the Fe$_{100-x}$Co$_x$ alloys (up to $x=60\%$) can hold an ultra-low intrinsic damping, of the order of $10^{-4}$, at $x\sim25\%$ \cite{schoen2016ultra}, despite being metallic. This surprising result, comparable to the values found in ferrimagnetic insulators (such as yttrium iron garnet \cite{Kajiwara2010}), was explained by the existence, at this particular value of $x$, of a sharp minimum density of states at the Fermi level, and corroborated by different methods and levels of theory \cite{PhysRevLett.107.066603, PhysRevB.87.014430, PhysRevB.98.104406, PhysRevB.108.014433, PhysRevB.98.174412, PhysRevB.92.214407}. In addition, Fe-Co has recently attracted attention with the observation of a giant anisotropy effect in $\alpha$ \cite{PhysRevLett.122.117203} in Fe$_{50}$Co$_{50}$ thin films, initially attributed to the variation of spin-orbit coupling due to local distortions in the alloy. This was also experimentally verified for other similar alloys involving Fe and Co \cite{PhysRevApplied.19.024030,Wang2023}, but the exact mechanisms responsible for this effect is still an open question \cite{PhysRevB.103.L220405,PhysRevB.103.L220403,PhysRevApplied.19.024030,PhysRevLett.131.186703}. Additionally, at the theoretical level, Lu \textit{et al.} \cite{PhysRevB.108.014433} suggested that for $x\in\{30\%,50\%\}$ the explicit treatment of damping as a nonlocal parameter in spin-dynamics have the effect of increased magnon lifetime predictions (for specific wave vectors $\mathbf{q}$). Together, those observations clearly consolidate FeCo as a benchmark for Gilbert damping investigations, encompassing many hidden features across the alloy series.

Despite the progress in studying these interesting phenomena in FeCo, other potential sources of hidden features are still not quite well-explored in the scope of Gilbert damping. Some have been recently highlighted as effects that are or might be relevant, also accounting for site-nonlocal damping  \cite{PhysRevB.109.094417,PhysRevB.103.L220405}. They are often related to changes in the local environment: reduced dimensionality, chemical disorder, and thermal fluctuations in the lattice (atomic displacements) and in the magnetic state. Although the former remains as the least discussed among the known effects \cite{PhysRevB.109.094417}, the other two, which constitute causes for electron scattering processes, are generally addressed by effective alloy analogy models (see, e.g., Refs. \cite{PhysRevB.91.165132,PhysRevLett.107.066603}). In particular, an explicit atomistic picture is still missing, and here we choose to address this knowledge gap in the context of the \textit{chemical disorder}. 

In this sense, the alloy models generally used -- the virtual crystal approximation \cite{PhysRevB.45.12911} or the coherent potential (CPA) approximation \cite{faulkner1982modern} -- were originally developed to describe chemical disorder in materials \cite{PhysRevLett.52.77} via an effective medium. The effective medium is a result of locally averaging possible clusters with different configurations and compositions within the alloy, and the CPA is designed to give the scattering properties of electron states that correspond to a proper configuration average of the alloy. This theoretical description 
works reasonably well to reproduce experimentally measured electronic structure properties of many alloys \cite{PhysRevLett.93.027203,PhysRevB.36.9657,PhysRevLett.40.900,PhysRevB.66.214206} and related materials properties, including Gilbert damping \cite{PhysRevB.87.014430,PhysRevB.92.214407,PhysRevB.94.214410,PhysRevB.108.014433}. Nevertheless, they frequently overlook the impacts of the local environment, including atom clustering, among others -- precisely the conditions under which phenomena like damping anisotropy were originally suggested to arise \cite{PhysRevLett.122.117203}. 

The local chemical environment has been shown to modify the magnetic properties of a given material, such as the atomic magnetic moment, magnetic exchange interactions \cite{bezerra2013complex, miranda2017dimensionality, de2016first, martins2016magnetic}, as well as  magnetic anisotropy \cite{PhysRevB.97.224402}. For instance, the presence of heterogeneous easy-axis orientations in clusters leads to a vanishing low magnetic anisotropy in permalloy (Ni$_{80}$Fe$_{20}$), although the favored easy-axis orientations at each site, dependent on the local atomic arrangement and the cluster's composition, can be several times larger than that of elemental Fe or Ni \cite{PhysRevB.97.224402}. As a result, the spin Hamiltonian and, thus, the energy landscape is modified, impacting the dynamics of the local magnetic moments in a system described by the LLG equation \cite{PhysRevB.86.224401}. In the context of FeCo, early calculations of ordered bcc Fe$_{50}$Co$_{50}$ structures have demonstrated that the magnetic anisotropy can change by some two orders of magnitude when particular Fe/Co configurations are considered \cite{Razee2001}. Also, experimental studies have shown that the magnetic properties are influenced by various local environment factors, including composition, morphology \cite{hesani2010effect, PhysRevLett.96.257205, karipoth2016magnetic}, as well as the size of Fe-Co alloy nanoclusters. 

In this study, we utilize Fe-Co alloys as a platform to investigate the impact of local compositions and atom arrangements on $\alpha$ via \textit{ab-initio} calculations. Here, we suggest the existence of hidden features (e.g., enhanced or decreased damping values) when an explicit atomistic picture is considered. As a consequence, local chemical engineering may be considered as a possible route to manipulate the Gilbert damping of such materials. 

The paper is organized as follows. In Sec. \ref{sec:methods}, we provide details on the density functional theory (DFT) calculations, the embedded cluster virtual crystal approximation model, along with the considered atomic configurations with different Co concentrations. Section \ref{sec:results} presents the results of the local environment effects on the Gilbert damping. In Sec. \ref{sec:conclusion}, we give a summary and an outlook.